\documentclass[10pt,a4paper]{amsart}
\usepackage[margin=19mm]{geometry}
\usepackage{amsmath,amssymb,mathtools}
\usepackage[scr=rsfs,cal=euler]{mathalfa}
\usepackage{xfrac}
\usepackage[unicode,hidelinks,bookmarks=false]{hyperref}
\newcommand{\ie}{i.e.\ }
\newcommand{\cf}{cf.\ }
\newcommand{\resp}{respectively\ }
\newcommand{\Subsection}{\subsection}
\providecommand{\nobreakdash}{\nobreak\hbox{-}}
\setlength{\emergencystretch}{2em}
\multlinegap0pt
\usepackage{mdframed}
\usepackage{mdframed}
\usepackage{mdframed}
\usepackage{mdframed}
\usepackage{amssymb}
\usepackage{mathtools}
\usepackage{tikz}
\usepackage{mdframed}
\usepackage{tcolorbox}
\usepackage{xcolor}
\usepackage[shortlabels]{enumitem}
\newtheorem{claim}{Claim}
\newtheorem{proofclaim}{Proofclaim}
\title{On the number of distinct spanning trees in pseudorandom graphs}
\author{Yiting Wang}
\date{Source: arXiv:2605.12742v1 (CC BY 4.0)}
\begin{document}
\maketitle

\noindent\emph{Source and licence.} On the number of distinct spanning trees in pseudorandom graphs. Version arXiv:2605.12742v1 (CC BY 4.0), distributed by arXiv under the Creative Commons Attribution 4.0 International licence, http://creativecommons.org/licenses/by/4.0/, as recorded in the arXiv licence metadata for this exact version and shown on the abstract page https://arxiv.org/abs/2605.12742v1. Full licence notice: \texttt{licenses/arxiv-2605.12742-CC-BY-4.0.txt}.\par\smallskip

\noindent\textit{Editorial scope.} Counted unit: the article's claim claim (source0.tex lines 266--268). The article's complete original proof of it is delivered with it (source0.tex lines 269--272, one \texttt{proof} environment of 4 lines). No mathematics is written, repaired or re-derived: the statement and the proof are the article's own lines, in the article's own notation, and no retained line is substituted or re-flowed.  No further source-local context is needed: every cross-reference of every retained line resolves inside the two delivered spans.  Nothing was omitted from any retained span, so the package carries no omission record.  No retained line contains a citation, so the wrapper carries no bibliography and no bibliography entry.  No retained line contains a figure environment, an image-inclusion command or drawing code, so no image is delivered and the package declares no asset dependency.  Editorial formatting only, in addition: an AMS \texttt{amsart} preamble that restates the article's own declarations and macros used by the retained lines, namely the environments \texttt{claim}, \texttt{proofclaim} on separate counters, together with the standard packages those lines need.  The retained environments print in this document as Claim 1, Proofclaim 1 in order of appearance, whereas the article prints the same items with the numbering of its own full document.  The complete native source of the article as distributed in the arXiv e-print for this exact version is bundled unchanged as \texttt{sources/200443/source0.tex}.
    \begin{claim}
        The map $(T_1,\dots, T_L)\mapsto T$ is injective.
    \end{claim}

    \begin{proofclaim}
        First, since $T_i$ has only $K$ vertices, the only maximal pendant paths of length greater than $K$ in $T$ are the two paths added in Step 3. Hence these two paths and their attachment vertices $\{x_1,x_L\}$ are identifiable from $T$.
        Moreover, since $b = R-a \geq 4K-1> a$, the two paths have distinct lengths. From there, we can distinguish $x_1$ from $x_L$. Lastly, the spine is the unique path connecting $x_1$ to $x_L$ in $T$ and so the ordered sequence $x_1,x_2,\dots, x_L$ is identifiable. Consequently, every distinct choice of $(T_1,\dots, T_L)$ gives rise to a distinct tree $T$, which proves the map is injective.
    \end{proofclaim}

\end{document}
